% ===========================================================================
\section{Performance Analysis}
% Source material: docs/evaluation-parameters.md, bench/results/
% This is the section the final evaluation weighs most heavily. Every number
% needs a stated environment and a stated sample size.
% ===========================================================================

\subsection{Setup}

\stub{Describe the measurement environment precisely enough to be reproducible:
  hardware, client versions, network, RPC provider, gas price at the time of
  measurement, contract commit hash, sample sizes, and how repetitions were
  aggregated.}

\begin{table}[htbp]
  \centering
  \caption{Measurement environments.}
  \label{tab:bench-setup}
  \small
  \begin{tabular}{@{}L{3.7cm}L{5.3cm}L{6.9cm}@{}}
    \toprule
    \textbf{Environment} & \textbf{Used for} & \textbf{Configuration}                                     \\
    \midrule
    Local EVM            & Gas measurements  & \note{toolchain version, solc version, optimizer settings} \\
    Local node           & Throughput        & \note{host specification}                                  \\
    Public testbed       & Latency, cost     & \note{network, RPC provider, gas price, block time}        \\
    \bottomrule
  \end{tabular}
\end{table}

\stub{State plainly which numbers come from a local environment and which from a
  live network, and why the split was made that way. Do not present a locally
  measured number as though it came from the testbed.}

\subsection{Results by evaluation parameter}

Results are presented against the parameters defined in the Evaluation 1 report.
Targets stated there were hypotheses; where a target was missed, the result is
reported as measured and analysed.

\subsubsection{Cost (A1)}

\stub{Report the gas per operation (A1). A fiat cost per vote is an
  illustration at a stated gas price, not a graded figure, since it depends on a
  gas price and an exchange rate outside the project's control.}

\subsubsection{Performance (B1, B2)}

\stub{Report latency distributions rather than averages: median and tail. Give
  proof generation time by device class. Report the throughput curve and say where
  it bends.}

\subsubsection{Scalability (C1--C3)}

This is the central claim of the design.

\begin{figure}[htbp]
  \centering
  \fbox{\parbox[c][5.5cm][c]{0.9\linewidth}{\centering\itshape
      Figure placeholder: gas cost against voter-roll size.\\
      Generate from \texttt{bench/results/} by script, not by hand.}}
  \caption{Gas cost per operation against the size of the electorate.}
  \label{fig:c1}
\end{figure}

\stub{State whether gas is flat in the size of the electorate, and give the
  measured figures at each roll size. If the line is not flat, explain precisely
  where the growth enters and what it would take to remove it. A negative result
  reported clearly is a better outcome than a favourable one that was not measured.}

\subsubsection{Security and correctness (D1--D5)}

\stub{For each property, state how it was demonstrated and what the outcome was.
  Where a property is only partially achieved, say so and give the limit precisely.}

\subsubsection{Usability (E1--E3)}

\stub{Report the user-testing results, including the tasks participants failed
  and what caused each failure. Failures here are the most useful data in the
  section.}

\subsection{Discussion}

\stub{Step back from the individual numbers. Which design decisions were
  vindicated by measurement, which were not, and what would be done differently
  with the benefit of the results? Name the system's principal remaining
  limitation.}
